\section{Conclusion}
\label{sec:conclusion}

Intrinsic-reward shaping can substantially accelerate learning when extrinsic feedback is sparse or delayed, but simple novelty- or prediction-error-based bonuses can also concentrate exploration in regions where progress is slow, unstable, or effectively noisy.
We introduced the \emph{GLPE family}, which ties intrinsic reward to two complementary signals computed from on-policy transitions: feature-space impact and region-local learning progress estimated from short- and long-horizon prediction-error averages in an online latent-space partition.
The full GLPE variant adds a simple region-wise gate that suppresses intrinsic shaping when prediction error stays high while learning progress stalls, providing a targeted guardrail against unproductive curiosity; GLPE (no gate) uses the same base score without this additional thresholding.

Across five Gymnasium control benchmarks and a shared PPO backbone, GLPE (no gate) is the most consistent default, remaining competitive with strong intrinsic baselines under both step-normalized and wall-clock-normalized views.
The gated variant is most valuable in sparse-reward regimes such as \textsc{MountainCar-v0}, where it improves reliability and reduces instability relative to ungated shaping, but it can be unnecessarily conservative in high-variance settings such as \textsc{Humanoid-v5}.
Finally, the cosine taper schedule makes intrinsic shaping intentionally transient, smoothly reducing its contribution late in training so that final performance reflects primarily the learned task policy.

This study is limited to vector-observation benchmarks and on-policy PPO, and (as with most intrinsic bonuses) the shaping signal we apply is not policy-invariant.
Future work can extend GLPE to richer observation modalities and alternative training regimes, and can make the gated variant more computationally efficient and adaptive without changing the underlying intrinsic objective.
